\section{Homogeneous Samples are Optimal for Small-Budget Pilots}
\label{sec:small-budget}

The key insight we derive in our model is that whether representative sampling is optimal depends critically on how resource-constrained the pilot study is.
While in the large-budget regime representative sampling is near-optimal, in this section we show that for sufficiently small pilot budgets, the optimal pilot design concentrates sampling on a single type.

We begin in \Cref{sec:small-budget-optimal-design} by deriving the optimal small-budget design when an RCT is used to determine whether to receive \emph{any} unknown downstream payoff under a general prior distribution with moment assumptions.
In \Cref{sec:small-budget-optimal-pilot} we specialize the result to give the optimal design for maximizing the pilot RCT impact in our model, where the payoff is the result of running a follow-up RCT and making an adoption decision based on its statistical significance.
We then discuss in \Cref{sec:small-budget-heterogeneity} how the optimal pilot design depends on the prior heterogeneity of the treatment effects under additional assumptions on the prior distribution.

\subsection{Small-budget optimal design for general downstream payoffs}
\label{sec:small-budget-optimal-design}

Consider an RCT sample $\bm{s}_1$ with total sample size $n_1=\bm{1}^\top\bm{s}_1$ and unknown vector of heterogeneous treatment effects $\bm{\tau}$.
As described in \Cref{sec:Model}, consider a statistical test of the form
\[
  \Test_1 \defeq \indic{\sqrt{n_1}\cdot \widehat{\ATE}_1\ge z_{1-\alpha_1}}.
\]
However, assume the test determines whether the decision maker receives a general downstream payoff $Y$, and that the objective is to maximize the expected realized payoff,
\[
  V_Y(\bm{s}_1)
  \defeq
  \E\left[
    \Test_1\cdot Y
  \right].
\]
We derive the optimal small-budget design for this objective, under the assumption that the payoff $Y$ is \emph{non-adaptive}, i.e. it does not depend directly on the RCT observations or, more formally, the payoff is independent of the pilot outcomes conditioned on the treatment effect vector $\bm{\tau}$.

We begin with a heuristic derivation.
The test statistic can be written as a linear function of the RCT sample divided by the square root of the total sample size, plus standard normal noise
\[
  \sqrt{n_1}\cdot\widehat{\ATE}_1
  =
  \sqrt{n_1}\cdot\frac{\bm{s}_1^\top\bm\tau}{n_1}+\varepsilon_1
  =
  \frac{\bm{\tau}^\top\bm{s}_1}{\sqrt{\bm{1}^\top\bm{s}_1}}+\varepsilon_1,
  \qquad
  \varepsilon_1\mid \bm\tau\sim\mathcal{N}(0,1).
\]
When samples are small, so is the expected value---the non-centrality---of this test statistic.
A Taylor expansion of the objective around zero non-centrality shows that in the small-budget regime, the objective is approximately equivalent to maximizing the expected product of the test statistic non-centrality and the downstream payoff.
\begin{alignat*}{2}
  V_Y(\bm{s}_1)
  &=
  \E\left[\indic{\frac{\bm{\tau}^\top\bm{s}_1}{\sqrt{\bm{1}^\top\bm{s}_1}}+\varepsilon_1 \ge z_{1-\alpha_1}} \cdot Y\right]
  &\qquad&
  (\text{definition of }\Test_1) \\
  &=
  \E\left[
    \Phi\!\left(
      \frac{\bm{\tau}^\top\bm{s}_1}{\sqrt{\bm{1}^\top\bm{s}_1}}
      - z_{1-\alpha_1}
    \right)
    \cdot Y
  \right]
  &\qquad&
  (\text{non-adaptivity, }Y\ind \varepsilon_1 \vert \bm{\tau}) \\
  &\approx
  \alpha_1\E[Y]
  +\varphi(z_{1-\alpha_1})\,
  \E\left[
    \frac{\bm{\tau}^\top\bm{s}_1}{\sqrt{\bm{1}^\top\bm{s}_1}}
    \cdot Y
  \right]
  &\qquad&
  (\text{expansion at zero non-centrality}) \\
  &=
  \alpha_1\E[Y]
  +\varphi(z_{1-\alpha_1})\,
  \frac{\E[Y \bm\tau]^\top\bm{s}_1}{\sqrt{\bm{1}^\top\bm{s}_1}}
  &\qquad&
  (\text{linearity of expectation})
\end{alignat*}

This means that for small budgets, maximizing the expected payoff is approximately equivalent to maximizing a linear signal divided by the square root of the sample size:
\begin{align*}
  \underset{\bm{s}_1}{\maximize}
  \quad
  \frac{\E[Y \bm\tau]^\top\bm{s}_1}{\sqrt{\bm{1}^\top\bm{s}_1}}
  \qquad
  \subjectto
  \quad
  \bm c^\top\bm{s}_1\le b,
  \quad
  \bm{s}_1\in\R_{\geq 0}^{D}\setminus\{\bm{0}\}.
\end{align*}
The ratio of a linear function to the square root of a linear function is quasiconvex when both linear functions are positive, and quasiconvex functions are maximized over polytopes at extreme points.
Thus, so long as there is at least one type with positive signal $\E[Y \tau_d]>0$, the approximate objective is maximized by investing the entire budget into sampling a single type: the one that maximizes the \emph{small-budget index}
\[
  \pi_d^Y
  \defeq
  \frac{\E[Y \tau_d]}{\sqrt{c_{d}}}.
\]

Formalizing the argument above requires care on two points: showing that the objective does not approach a maximizer as the sample approaches zero $\bm{s}_1\to 0$ where the objective is undefined, and establishing uniform control of the remainder of the Taylor expansion to show that the optimizer of the objective and its leading-order term are exactly equal for sufficiently small finite budgets.
We do so by writing the objective in terms of the investment $\invst=\bm{c}^\top \bm{s}_1$ and the unit-investment normalized sample composition $\overline{\bm{s}}_1$ as described in \Cref{sec:Model}, and taking a Taylor expansion in terms of $\sqrt{I}$.

\begin{proposition}[Small-investment expansion for non-adaptive payoffs]
  \label{prop:small-investment-expansion}
  Fix a prior distribution of $\bm\tau$, and let $Y$ be non-adaptive with $\E[|Y|]<\infty$ and $\E[|Y|\|\bm\tau\|_2^3]<\infty$.
  Then there exists an investment threshold $\overline{\invst}>0$ such that for all investments $0 < \invst \le \overline{\invst}$ and all unit-investment samples $\overline{\bm{s}}_1\in\overline{\mathcal{S}}$,
  \[
    V_Y(\invst\,\overline{\bm{s}}_1)
    =
    \alpha_1\, \E[Y]
    +\sqrt{\invst}\,\varphi(z_{1-\alpha_1})\,
    \frac{\E[Y\bm\tau]^{\!\top}\overline{\bm{s}}_1}{\sqrt{\bm{1}^{\top}\overline{\bm{s}}_1}}
    +R(\invst,\overline{\bm{s}}_1),
  \]
  where $R(\invst,\overline{\bm{s}}_1)$ is $O(\invst)$ and Lipschitz in $\invst$ (both uniformly in $\overline{\bm{s}}_1$) and $O(\invst)$-Lipschitz in $\overline{\bm{s}}_1$.
\end{proposition}

This result and the following small budget optimal policy are proved in \Cref{\sectionof{prop:small-investment-expansion-appendix}}.
The expansion of the expected payoff objective allows us to show that for all budgets~$b>0$ below a fixed budget threshold~$b<\overline{b}$ depending on the problem parameters, the optimal design is completely determined by the first non-constant term in the expansion,
\[
  \sqrt{\invst}\,\cdot\,
  \frac{\E[Y\bm\tau]^{\!\top}\overline{\bm{s}}_1}{\sqrt{\bm{1}^{\top}\overline{\bm{s}}_1}}.
\]
Similarly to our heuristic derivation above, we argue that the linear over square-root-of-linear term is maximized at an extreme point.
Here, however, it is a function on the set of unit-investment designs $\overline{\mathcal{S}}=\{\bm{x}\in\R^D_{\ge 0} : \bm{c}^{\top}\bm{x}=1\}$ which is a compact polytope and does not contain the zero sample for which our objective is undefined.
The term scales with the square root of the investment, and is therefore maximized by investing the entire budget.
Thus, the optimal design invests the entire budget in a single type, and the Lipschitz properties of the remainder allow us to extend from the term above to the original expected payoff objective under a sufficiently small budget.
This brings us to the general small-budget result.

\begin{theorem}[Optimal small-budget design for non-adaptive payoffs]
  \label{thm:optimal-small-budget-design-general}
  Fix a prior distribution of $\bm\tau$, and let $Y$ satisfy the assumptions of \Cref{prop:small-investment-expansion}.
  Suppose there exists a unique type $d^\star$ maximizing the \emph{small-budget index}
  \[
    \pi_d^Y
    \defeq
    \frac{\E[Y \tau_d]}{\sqrt{c_{d}}}
    \in
    \R.
  \]
  Suppose also that $\pi_{d^\star}^Y>0$.
  Then there exists a budget threshold $\overline{b}>0$ such that for all budgets $0 < b \le \overline{b}$, a unique optimal design exists for $V_Y$ and invests the entire budget in the type $d^\star$.
\end{theorem}

\subsection{Small-budget optimal pilot}
\label{sec:small-budget-optimal-pilot}

For the pilot impact objective, the downstream payoff is $Y=\Test_2\cdot\ATE_3$, so the small-budget index becomes
\[
  \idx_d
  \defeq
  \frac{\E[\tau_d \cdot \Test_2 \cdot \ATE_3]}{\sqrt{c_{d}}}
  \in
  \R.
\]

\begin{corollary}[Optimal small-budget pilot]
  \label{cor:optimal-small-budget-pilot}
  Suppose the prior distribution of $\bm\tau$ has finite fourth moments, and that there exists a unique type $d^\star$ maximizing the small-budget index $\idx_d$.
  Then there exists a budget threshold $\overline{b}>0$ such that for all pilot budgets $b<\overline{b}$, a unique optimal pilot exists and invests the entire budget in the type $d^\star$.
\end{corollary}

This follows from \Cref{thm:optimal-small-budget-design-general}.
The probability model and finite fourth moments imply the non-adaptivity and moment assumptions for $Y=\Test_2\cdot\ATE_3$, and the positivity condition is automatic for the pilot impact objective; see \Cref{\sectionof{prop:positive-impact-objectives-appendix}}.

Not only are we able to show that representative sampling is never optimal for pilots that screen treatments for follow-up study under a sufficiently small budget, in fact the other extreme is optimal here: sampling only from a subpopulation containing a single type.
This is counter to the many calls for the use of more representative samples in the social and health sciences~\cite{FDAHH20, DepartmentofEducationBo22,BryanTi21}, and reflects that modeling decision outcomes under a decision rule requiring statistical significance yields qualitatively different optimal designs than those usually found in work focused on externally valid inference~\cite{Tipton13,EgamiLe23,Bouyamourn25} and minimax-regret objectives~\cite{HuZh24,OleaPr25}.

Further, we are able to characterize exactly which type to sample.
Intuitively, large $\tau_d$ makes passing the pilot more likely, and we want this to happen when the target effect $\ATE_3$ is large under those realizations where the follow-up statistical test $\Test_2$ passes.
Thus, we trade off the expected value of $\tau_d \cdot \Test_2 \cdot \ATE_3$ against the square root of the sampling cost $c_d$, since the impact improves with the square root of the investment.
We now proceed to give a more interpretable characterization of the optimal type to sample under additional assumptions on the prior distribution of the treatment effects.

\subsection{The influence of treatment effect heterogeneity on
optimal pilot design}
\label{sec:small-budget-heterogeneity}

In addition to characterizing the optimal sampling type for the pilot, \Cref{cor:optimal-small-budget-pilot} also allows us to interpret how the optimal type to sample in the small-budget regime should depend on our treatment effect prior.
In the remainder of this section we analyze the form of $\pi_d$ by restricting the prior to the class of elliptical distributions~\cite{FangKo90}
\[
  \bm{\tau}\sim\mathrm{EC}_D(\bm{\mu},\Sigma,\phi),
\]
which includes and can be viewed as generalizing the multivariate normal distribution.
We begin by informally describing the dependence under typical experimentation conditions, before giving in \Cref{prop:marginal-impact-interpretation} a decomposition of the small-budget index from which this interpretation is derived.

To begin, assume the follow-up RCT uses a representative sample, captured in our model by assuming $\bm{s}_2\propto \bm{s}_3$.
Further, assume the prior expected target effect is not harmful, $\E[\ATE_3]\ge 0$, and that the target-population effect is nondegenerate, $\var(\ATE_3)>0$.

Heterogeneity in the conditional average treatment effects $\tau_d$ across types can come into the prior in two ways: \emph{expected} variation in the prior mean $\E[\tau_d]$ across types, and \emph{unknown} variation through the covariance matrix $\cov(\bm\tau)$.
Under the scenario described above, the small-budget index is proportional to a weighted sum of a type's prior mean and how predictive its treatment effect is of the target population effect, scaled by the square root of the sampling cost: for some positive weight $\overline W_1$,
\vspace{0.3cm}
\begin{equation}
  \idx_d \propto \frac{\overbrace{\overline W_1 \cdot \E[\tau_d]}^{\text{Expected
    heterogeneity}} \ + \
    \overbrace{\cov(\tau_d, \ATE_3)}^{\text{Unknown heterogeneity}}
  }{\underbrace{\sqrt{c_d}}_{\text{Cost heterogeneity}}}.
  \label{eq:small-budget-index-interpretation-special-case}
\end{equation}
\vspace{0.3cm}

To see how the prior covariance $\cov(\bm\tau)$ affects the expression above, consider the further specialization to a unit-variance exchangeable prior covariance with common correlation $\rho\in[0,1]$,
\[
  \cov(\bm\tau) =
  \begin{pmatrix}
    1 & \rho & \cdots & \rho \\
    \rho & 1 & \cdots & \rho \\
    \vdots & \vdots & \ddots & \vdots \\
    \rho & \rho & \cdots & 1
  \end{pmatrix}.
\]
In this case each type effect $\tau_d$ is predictive of the target population effect $\ATE_3$ by an amount depending on its share of the target population $s_{3d}/n_3$, attenuated by the common correlation $\rho$.
\begin{equation}
  \cov(\tau_d, \ATE_3)
  =
  \sum_{d'=1}^D \frac{s_{3d'}}{n_3}\cov(\tau_d, \tau_{d'})
  =
  \rho + (1-\rho)\frac{s_{3d}}{n_3}
  \label{eq:small-budget-index-interpretation-special-case-cov}
\end{equation}
These assumptions allow us to interpret the small-budget index under different kinds of elliptical priors:
\begin{enumerate}
  \item {
      When we do not expect much heterogeneity in the treatment effect across types, the prior means $\mu_d$ are roughly equal across types and the prior correlation $\rho$ is close to 1.
      In this case the numerator in \eqref{eq:small-budget-index-interpretation-special-case} is constant and the small-budget optimal policy depends only on cost: it reduces to sampling the cheapest type, i.e. convenience sampling.
    }
  \item {
      If we have prior knowledge that there is potential heterogeneity in the treatment effect but do not anticipate that it favors any specific type, the prior correlation $\rho$ is small and the prior means $\E[\tau_d]$ are still equal.
      In this case the optimal small-budget pilot skews towards types that make up a larger share of the target population.
      See \Cref{fig:exchangeable-2types-by-rho-zero-mean} for an example.
    }
  \item {
      Finally, if we have prior knowledge that the expected effect is larger in specific types than others, so $\E[\tau_d]$ is not equal across types, and we believe that the expected effect is not harmful $\E[\ATE_3]\ge 0$, the optimal small-budget pilot skews towards types with larger prior means.
      See \Cref{fig:exchangeable-2types-by-mu} for an example.
    }
\end{enumerate}
\begin{figure}[ht]
  \centering
  \begin{subfigure}[t]{0.48\textwidth}
    \centering
    \includegraphics[width=\linewidth]{figures/exchangeable_2types_by_rho_zero_mean.pdf}
    \caption{
      Decreasing $\rho$ (increasing heterogeneity) favors the type that has a larger target population share and is therefore more predictive.
      Here under $\bm\mu=\bm{0}$.
    }
    \label{fig:exchangeable-2types-by-rho-zero-mean}
  \end{subfigure}
  \hfill
  \begin{subfigure}[t]{0.48\textwidth}
    \centering
    \includegraphics[width=\linewidth]{figures/exchangeable_2types_by_mu2.pdf}
    \caption{
      Increasing $\mu_2$ (the prior mean of type 2) while keeping $\mu_1 = 0$ shifts the index toward type 2, even though type 1 is more common.
      Here under $\rho=0.2$.
    }
    \label{fig:exchangeable-2types-by-mu}
  \end{subfigure}
  \caption{
    The small-budget index for two types under a large ($n_2 = 100$), representative follow-up and a multivariate normal prior with unit prior standard deviations.
    For any sufficiently small budget, the optimal pilot invests the entire budget in sampling the type with the larger index.
  }
  \label{fig:exchangeable-2types}
\end{figure}

\Needspace{12\baselineskip}
The preceding interpretations are formalized in the following general decomposition of the small-budget index, proved in \Cref{\sectionof{prop:marginal-impact-interpretation-appendix}}.

\begin{proposition}[The small-budget index under an elliptical prior]
  \label{prop:marginal-impact-interpretation}
  Suppose the treatment effect prior is an elliptical distribution,
  \[
    \bm\tau\sim\mathrm{EC}_D(\bm\mu,\Sigma,\phi),
  \]
  with finite second moments.
  Then there exist scalars
  \(W_1,W_2,W_3\in\R\), not depending on \(d\), such that for every
  type \(d\),
  \[
    \idx_d
    =
    \frac{
      W_1\,\E[\tau_d]
      +
      W_2\,\cov(\tau_d,\ATE_2)
      +
      W_3\,\cov(\tau_d,\ATE_3)
    }{\sqrt{c_d}}.
  \]

  If additionally \(\bm s_2\propto \bm s_3\), \(\E[\ATE_3]\ge0\), and \(\var(\ATE_3)>0\), then the decomposition reduces to
  \[
    \idx_d
    \propto
    \frac{
      \overline W_1\,\E[\tau_d]
      +
      \cov(\tau_d,\ATE_3)
    }{\sqrt{c_d}},\qquad \overline W_1 > 0
  \]
  with a positive constant of proportionality, where the weight on the prior mean is strictly positive.
\end{proposition}

This decomposition shows that the small-budget index $\idx_d$ is the cost-adjusted, weighted sum of three interpretable components: the prior mean $\E[\tau_d]$ of the type's treatment effect, how predictive it is of the follow-up average effect $\ATE_2$, and how predictive it is of the target population effect $\ATE_3$.
Furthermore, when the follow-up is representative and the treatment is not expected to harm the target population, the index prioritizes types with larger prior means.
